\section{�bungsaufgaben}\label{kap:aero_uebungen}

\begin{prob} %�bung1
\einfach{aero_testpotential}{�berpr�fen auf die Existenz eines Potentials}

In Beispiel 2.4 wurde f�r die station�re Str�mung zwischen zwei W�nden
im Abstand $b$ das Geschwindigkeitsfeld
\begin{equation*}
  v_z(y) = \frac{1}{2\eta} \frac{\partial p(z)}{\partial z}
  \lk \lk y - \frac{b}{2} \rk^2 -  \frac{b^2}{4} \rk
\end{equation*}
gefunden. Analog ergab sich f�r ein Kreisrundes Rohr mit dem Durchmesser $R$
das Feld
\begin{equation*}
  v_z(r) = \frac{1}{4\eta} \frac{\partial p(z)}{\partial z}
  \lk r^2 - R^2 \rk \; .
\end{equation*}
�berpr�fen Sie beide Geschwindigkeitsfelder auf die Existens eines Potentials
und geben Sie dieses gegebenenfalls an.

\hyperref[lsg:aero_testpotential]{\textcolor{blue} {$\rightarrow$ L�sungsvorschlag}}\\
%\href{FingerUeb3_AB.pdf#lsg:aero_testpotential}{\textcolor{blue} {$\rightarrow$ L�sungsteil}}.\\

\noindent\rule[1ex]{\textwidth}{0.5pt}
\end{prob}

%%%%%%%%%%%%%%%%%%%%%%%%%%%%%%%%%%%%%%%%%%%%%%%%%%%%%%%%%%%%%%%%%%%%%%%%%%%%%


\begin{prob} %�bung2
\einfach{aero_potentialbeispiel}{Einfache Beispiele f�r Potentialstr�munen}

Stellen Sie die Potentialgleichungen f�r die Beispiele der Str�mung zwischen
zwei Platten und durch ein Kreisrundes Rohr aus \probref{aero_testpotential}
auf. �berzeugen Sie sich dabei am konkreten Beispiel, dass der Reibhungsterm
verschwindet. L�sen Sie die Potentialgleichung und berechnen Sie daraus
das Geschwindigkeitsfeld.

L�sen Sie anschlie�end die Navier-Stokes-Gleichungen ohne Reibhungsterm (also
die zugeh�rigen Euler-Gleichungen) und zeigen Sie, dass sie auf dieselbe
L�sung kommen.

\hyperref[lsg:aero_potentialbeispiel]{\textcolor{blue} {$\rightarrow$ L�sungsvorschlag}}\\
%\href{FingerUeb3_AB.pdf#lsg:aero_potentialbeispiel}{\textcolor{blue} {$\rightarrow$ L�sungsteil}}.\\

\noindent\rule[1ex]{\textwidth}{0.5pt}
\end{prob}

%%%%%%%%%%%%%%%%%%%%%%%%%%%%%%%%%%%%%%%%%%%%%%%%%%%%%%%%%%%%%%%%%%%%%%%%%%%%%

\begin{prob}%�bung3
\einfach{aero_haftwand}{Str�mung und Haftbedingung an einer Wand}

L�sen Sie die Navier-Stokes-Gleichungen f�r den ebenen Fall numerisch, sodass
Sie das Verhalten einer Str�mung an einer Wand darstellen k�nnen. Setzen Sie
die korrekte Haftbedingung an einer Wand an und gehen Sie von den
Randbedingungen aus, die sich aus Beispiel~\ref{bsp:aero_HaftbedingungWand}
ergeben.

\hyperref[lsg:aero_haftwand]{\textcolor{blue} {$\rightarrow$ L�sungsvorschlag}}\\
%\href{FingerUeb3_AB.pdf#lsg:aero_haftwand}{\textcolor{blue} {$\rightarrow$ L�sungsteil}}.\\

\noindent\rule[1ex]{\textwidth}{0.5pt}
\end{prob}
